\section{Further questions}
The transition counting error obtained here is
$O((\log\tau)^{2/3})$ on fixed compact subsets of the outer bands.
Determining the sharper error for this signed operator requires more than
the smooth comparison used in Theorem~\ref{v7:thm:cumulative-count}.
Uniform laws for thresholds approaching zero or a band edge remain open.
The finite-frequency results determine the bulk graph and low-order
logarithmic terms; a full second-order law and transition classification
for an arbitrary commensurate family are not asserted.

\subsection{Boundary questions}
The isolated phase-zero mode has a certified nonzero amplitude; the
corresponding statement at arbitrary phase or for every gap eigenvalue
is not proved. Determining the absolute boundary indices would identify
specific fixed-offset ranks near the central transition. Uniform laws
as the spectral parameter approaches zero or an outer band require
additional analysis.
Other questions concern an elementary evaluation and convergence rate
for the phase-dependent Fredholm product, and extension of the
reflected limit to the full band-local Lipschitz class.
The kernel also has the exact integrable representation
\[
 K_\tau(x,y)=\frac{\mathbf u_\tau(x)^{\mathsf T}\mathbf v_\tau(y)}{x-y},
 \quad
 \mathbf u_\tau(x)=\begin{pmatrix}\sin(2\pi\tau x)\\1\end{pmatrix},
 \quad
 \mathbf v_\tau(x)=\frac1\pi\begin{pmatrix}1\\-\sin(2\pi\tau x)\end{pmatrix}.
\]
Its diagonal numerator vanishes, as required for a rank-two integrable
kernel. Evaluating the global connection in the boundary equation
of Theorem~\ref{v12:thm:explicit-gap-equation} remains open.
The scalar differential obstruction in
Section~\ref{boundary:app:commutator} leaves matrix, higher-order
and nonlocal reductions available. Generalized sine-kernel
Riemann--Hilbert methods \cite{KitanineEtAl2009} provide a methodological
comparison, not an evaluated connection formula for this boundary problem.



